\documentclass[12pt,a4paper]{article}
\usepackage{mathwizards-formulario}
\title{\textbf{Integrales Impropias — Ejercicios Resueltos}}
\author{}
\date{}

% ── Comandos propios de este documento ─────────
\newcommand{\converge}[1]{%
  \textcolor{resultcolor}{\textbf{Converge.} Resultado: $#1$}%
}
\newcommand{\diverge}{%
  \textcolor{divcolor}{\textbf{Diverge.}}%
}

\begin{document}
\maketitle

%---------------------------------------------------------------------------
% EJERCICIO 1
%---------------------------------------------------------------------------
\section*{Ejercicio 1}
\[
  \int_{1}^{+\infty} \frac{x^3}{x^5}\,dx
\]

\textbf{Paso 1 — Simplificar el integrando.}
\[
  \frac{x^3}{x^5} = x^{-2}
\]

\textbf{Paso 2 — Escribir como límite (límite en $+\infty$).}
\[
  \lim_{b\to+\infty} \int_{1}^{b} x^{-2}\,dx
\]

\textbf{Paso 3 — Integrar.}
\[
  \int x^{-2}\,dx = \frac{x^{-1}}{-1} = -\frac{1}{x}
\]

\textbf{Paso 4 — Evaluar y tomar el límite.}
\[
  \lim_{b\to+\infty} \left[-\frac{1}{x}\right]_{1}^{b}
  = \lim_{b\to+\infty} \left(-\frac{1}{b} + 1\right)
  = 0 + 1 = 1
\]

\textbf{Regla $p$:} es una integral $\int_1^{+\infty} x^{-p}dx$ con $p=2>1$ $\Rightarrow$ converge.

\medskip\converge{1}

\bigskip
%---------------------------------------------------------------------------
% EJERCICIO 2
%---------------------------------------------------------------------------
\section*{Ejercicio 2}
\[
  \int_{0}^{+\infty} \frac{x^3}{2x^7}\,dx
\]

\textbf{Paso 1 — Simplificar.}
\[
  \frac{x^3}{2x^7} = \frac{1}{2}\,x^{-4}
\]

\textbf{Paso 2 — Detectar singularidades.}
El integrando $\frac{1}{2}x^{-4}$ tiene una discontinuidad en $x=0$. La integral es impropia \emph{tanto} en el límite inferior ($x\to 0^+$) como en el límite superior ($+\infty$). Se divide:
\[
  \int_{0}^{+\infty} \frac{1}{2}x^{-4}\,dx
  = \int_{0}^{1} \frac{1}{2}x^{-4}\,dx + \int_{1}^{+\infty} \frac{1}{2}x^{-4}\,dx
\]

\textbf{Paso 3 — Analizar la parte en $x\to 0^+$.}
\[
  \lim_{\varepsilon\to 0^+}\int_{\varepsilon}^{1} \frac{1}{2}x^{-4}\,dx
  = \lim_{\varepsilon\to 0^+} \frac{1}{2}\left[\frac{x^{-3}}{-3}\right]_{\varepsilon}^{1}
  = \lim_{\varepsilon\to 0^+} \frac{1}{2}\left(-\frac{1}{3} + \frac{1}{3\varepsilon^3}\right)
  = +\infty
\]

La primera parte ya diverge.

\medskip\diverge

\bigskip
%---------------------------------------------------------------------------
% EJERCICIO 3
%---------------------------------------------------------------------------
\section*{Ejercicio 3}
\[
  \int_{1}^{+\infty} x^4\,dx
\]

\textbf{Paso 1 — Escribir como límite.}
\[
  \lim_{b\to+\infty} \int_{1}^{b} x^4\,dx
\]

\textbf{Paso 2 — Integrar.}
\[
  \int x^4\,dx = \frac{x^5}{5}
\]

\textbf{Paso 3 — Evaluar.}
\[
  \lim_{b\to+\infty} \left[\frac{x^5}{5}\right]_{1}^{b}
  = \lim_{b\to+\infty}\frac{b^5 - 1}{5} = +\infty
\]

\textbf{Regla $p$:} integral de $x^p$ con $p=4>0$ crece sin cota $\Rightarrow$ diverge.

\medskip\diverge

\bigskip
%---------------------------------------------------------------------------
% EJERCICIO 4
%---------------------------------------------------------------------------
\section*{Ejercicio 4}
\[
  \int_{1}^{+\infty} (x^2 - 9)\,dx
\]

\textbf{Paso 1 — Escribir como límite.}
\[
  \lim_{b\to+\infty} \int_{1}^{b} (x^2 - 9)\,dx
\]

\textbf{Paso 2 — Integrar.}
\[
  \int (x^2-9)\,dx = \frac{x^3}{3} - 9x
\]

\textbf{Paso 3 — Evaluar.}
\[
  \lim_{b\to+\infty}\left[\frac{x^3}{3}-9x\right]_{1}^{b}
  = \lim_{b\to+\infty}\left(\frac{b^3}{3}-9b\right) - \left(\frac{1}{3}-9\right)
  = +\infty
\]

El término $b^3/3$ domina y crece sin cota.

\medskip\diverge

\bigskip
%---------------------------------------------------------------------------
% EJERCICIO 5
%---------------------------------------------------------------------------
\section*{Ejercicio 5}
\[
  \int_{-\infty}^{1} \frac{2x}{\sqrt{x}}\,dx
\]

\textbf{Paso 1 — Dominio del integrando.}
$\sqrt{x}$ requiere $x \geq 0$. Para $x < 0$ el integrando no está definido en $\mathbb{R}$.
La integral desde $-\infty$ cruza la región $(-\infty, 0)$ donde el integrando es indefinido, por lo que la integral \emph{no existe}.

\medskip\diverge

\bigskip
%---------------------------------------------------------------------------
% EJERCICIO 6
%---------------------------------------------------------------------------
\section*{Ejercicio 6}
\[
  \int_{-\infty}^{0} (x^{-2} + x^2)\,dx
\]

\textbf{Paso 1 — Detectar singularidades.}
$x^{-2} = 1/x^2$ es discontinua en $x=0$, que coincide con el límite superior. Se necesita:
\[
  \lim_{\varepsilon\to 0^-}\int_{-\infty}^{\varepsilon}(x^{-2}+x^2)\,dx
\]

\textbf{Paso 2 — Dividir y analizar la parte en $x\to 0^-$.}
\[
  \lim_{\varepsilon\to 0^-}\int_{-1}^{\varepsilon} x^{-2}\,dx
  = \lim_{\varepsilon\to 0^-}\left[-\frac{1}{x}\right]_{-1}^{\varepsilon}
  = \lim_{\varepsilon\to 0^-}\left(-\frac{1}{\varepsilon}-1\right)
\]
Como $\varepsilon\to 0^-$, $-1/\varepsilon \to +\infty$.

\medskip\diverge

\bigskip
%---------------------------------------------------------------------------
% EJERCICIO 7
%---------------------------------------------------------------------------
\section*{Ejercicio 7}
\[
  \int_{-\infty}^{0} \frac{\sqrt{4x^3}}{2x}\,dx
\]

\textbf{Paso 1 — Dominio.}
$\sqrt{4x^3} = 2\sqrt{x^3}$ requiere $x \geq 0$. Para $x < 0$ el radicando $x^3 < 0$ y la expresión no es real.

La integral abarca $(-\infty, 0)$ donde el integrando no está definido en $\mathbb{R}$.

\medskip\diverge

\bigskip
%---------------------------------------------------------------------------
% EJERCICIO 8
%---------------------------------------------------------------------------
\section*{Ejercicio 8}
\[
  \int_{-\infty}^{1} \frac{2x^3}{\sqrt{4x^4}}\,dx
\]

\textbf{Paso 1 — Simplificar.}
\[
  \sqrt{4x^4} = 2\sqrt{x^4} = 2x^2 \quad (x^2 \geq 0\ \forall x)
\]
\[
  \frac{2x^3}{2x^2} = x
\]
El integrando está definido para todo $x \in \mathbb{R}$.

\textbf{Paso 2 — Escribir como límite.}
\[
  \lim_{a\to-\infty}\int_{a}^{1} x\,dx
\]

\textbf{Paso 3 — Integrar y evaluar.}
\[
  \lim_{a\to-\infty}\left[\frac{x^2}{2}\right]_{a}^{1}
  = \lim_{a\to-\infty}\left(\frac{1}{2} - \frac{a^2}{2}\right)
  = -\infty
\]

\medskip\diverge

\bigskip
%---------------------------------------------------------------------------
% EJERCICIO 9
%---------------------------------------------------------------------------
\section*{Ejercicio 9}
\[
  \int_{-\infty}^{+\infty} \frac{1}{x^2}\,dx
\]

\textbf{Paso 1 — Separar en la discontinuidad $x=0$.}
\[
  \int_{-\infty}^{+\infty} \frac{1}{x^2}\,dx
  = \int_{-\infty}^{0} \frac{1}{x^2}\,dx + \int_{0}^{+\infty} \frac{1}{x^2}\,dx
\]
\emph{Ambas} partes presentan una discontinuidad en $x=0$.

\textbf{Paso 2 — Analizar $\int_{0}^{+\infty}$ (la parte derecha).}
\[
  \lim_{\varepsilon\to 0^+}\int_{\varepsilon}^{1}\frac{1}{x^2}\,dx
  = \lim_{\varepsilon\to 0^+}\left[-\frac{1}{x}\right]_{\varepsilon}^{1}
  = \lim_{\varepsilon\to 0^+}\left(-1+\frac{1}{\varepsilon}\right)
  = +\infty
\]

Una parte ya diverge; la integral total diverge.

\textbf{Nota:} no se puede aplicar el valor principal de Cauchy como si fuera una integral ordinaria de $-\infty$ a $+\infty$.

\medskip\diverge

\bigskip
%---------------------------------------------------------------------------
% EJERCICIO 10
%---------------------------------------------------------------------------
\section*{Ejercicio 10}
\[
  \int_{-\infty}^{+\infty} (x^2 - 2x)\,dx
\]

\textbf{Paso 1 — Separar.}
\[
  \int_{-\infty}^{+\infty}(x^2-2x)\,dx
  = \int_{-\infty}^{0}(x^2-2x)\,dx + \int_{0}^{+\infty}(x^2-2x)\,dx
\]

\textbf{Paso 2 — Analizar $\int_{0}^{+\infty}$.}
\[
  \lim_{b\to+\infty}\int_{0}^{b}(x^2-2x)\,dx
  = \lim_{b\to+\infty}\left[\frac{x^3}{3}-x^2\right]_{0}^{b}
  = \lim_{b\to+\infty}\left(\frac{b^3}{3}-b^2\right)
  = +\infty
\]

El término $b^3/3$ domina $\Rightarrow$ la parte derecha diverge.

\medskip\diverge

\end{document}
